\subsection{Height one and full first-level elliptic torsion}
The first layer already recovers the height-one character, while elliptic full torsion has additional constraints that depend on the particular base field.

\begin{proposition}[The height-one index]\label{prop:height-one-index}
Let $p$ be any prime and let $K$ satisfy the standing complete-field hypotheses, with $\mu_p\subset K$. For every height-one group $M$ in the allowed Barsotti--Tate filtration class, its special character on prime-to-$p$ inertia is determined by $M[p]$. Thus
\[
                         \nuMin(p,e_K,1)=0.
\]
This includes $p=2$.
\end{proposition}

\begin{theorem}[A full-three-torsion curve of defect twelve]\label{thm:explicit-hesse-twelve}
Put
\[
 K=\Q_3(i,\zeta),\qquad i^2=-1,\quad \zeta^2+\zeta+1=0,
 \qquad \pi=\zeta-1,\qquad t=1+i\pi.
\]
Then $e_K=2$ and the residue field is $\F_9$. The elliptic curve
\[
 E:\ y^2=x^3-27t(t^3+8)x+54(t^6-20t^3-8)
\]
has full rational $3$-torsion and additive potentially supersingular reduction of exact semistability defect twelve. The same exact defect holds after unramified base change.
\end{theorem}

\begin{theorem}[A variable Hesse defect]\label{thm:hesse-variable-defect}
Let $K$ be a complete discretely valued field of characteristic zero with algebraically closed residue field of characteristic three, containing $\mu_3$ and of absolute ramification index $e$. Choose a uniformizer $\varpi$ and an integer $r$ such that
\[
                         0<2r<e,\qquad 3\nmid e+r.
\]
Put $t=1+\varpi^r$. The elliptic curve
\[
 E_t:\quad y^2=x^3-27t(t^3+8)x+54(t^6-20t^3-8)
\]
has full rational $3$-torsion and potentially supersingular reduction of exact semistability defect
\[
                             m=\frac{12}{\gcd(4,e+r)}.
\]
There is an explicit good-reduction extension of this degree, and the ramification degree of every good-reduction extension is divisible by $m$.
\end{theorem}

\begin{corollary}[Every first-level defect in characteristic three]\label{cor:three-first-level-table}
Allow the complete field to vary, with arbitrary perfect residue field of characteristic three. An additive potentially good elliptic curve with full rational $3$-torsion and specified defect $m$ exists at exactly the indices in the following table:
\[
\begin{array}{c|c}
 m&\text{absolute ramification indices}\\\hline
 2&e\in2\Z,\ e\ge4\\
 3&e\in2\Z,\ e\ge4\\
 4&e\in2\Z,\ e\ge2\\
 6&e\in2\Z,\ e\ge2\\
 12&e\in2\Z,\ e\ge2.
\end{array}
\]
There are no other nontrivial defects. In particular every positive even index is attained. Over the particular field $\Q_3(\zeta_3)$ there is an example of defect six.
\end{corollary}

\begin{corollary}[The cyclotomic base field]\label{cor:cyclotomic-full-p}
For $p\ge5$, no elliptic curve defined over $\Q_p(\zeta_p)$ has both full rational $p$-torsion and additive potentially good reduction. For $p=3$ such a curve exists.
\end{corollary}

Section~\ref{sec:first-level-refinements} proves these statements. The Hesse family and its preservation of symplectic $3$-torsion structures are classical; see \citet[\S1, Theorem~1.1, p.~2]{AnemaTopTuijp2018}, over perfect fields of characteristic different from three. Here the generic field has characteristic zero. The proofs supply the points and the exact local good-reduction fields for the displayed curves. The first-level table is a direct specialization of the complete tame and wild existence criteria. The fixed cyclotomic field requires an additional residue-field argument at $p=5$: an existential example at the same absolute index over a different field would not settle that corollary. The classifications of \citet{OzekiYoshida2026} concern curves defined over $\Q_p$; they are neighboring context and are not used as a proof of this assertion for curves defined over $\Q_p(\zeta_p)$.
